% SOURCE_AUTHORITY: sga5_fr_workpass.tex, lines 15223--15233 (LF-normalized unit).
% SOURCE_SHA256: 791F4EFFC5E02832D5D77ED03518C8156D6F07E4C8238B03545DB93D883FBB28
Antes de demostrar estos lemas, indiquemos cómo se aplican a la demostración de la existencia de los sistemas $(K'_\nu)$ y $(h'_\nu)$. Se razona por inducción sobre $\nu$, suponiendo ya construidos $K'_0,h'_0,\ldots,K'_{\nu-1},h'_{\nu-1}$; se aplican los lemas con
\[
A=\mathbb Z/\ell^{\nu+1}\mathbb Z
\quad\text{y}\quad
I=\ell^\nu\mathbb Z/\ell^{\nu+1}\mathbb Z
\]
(de modo que $I^2=0$),
\[
A_0=\mathbb Z/\ell^\nu\mathbb Z.
\]
Tomemos primero $L=K_\nu$ y $M=K'_{\nu-1}$ en el primer lema. Teniendo en cuenta la observación 2, obtenemos la existencia de $K'_\nu=L'$ con el morfismo de transición $K'_\nu\to K'_{\nu-1}$ dado por $w$, de modo que se satisfaga la condición 1) enunciada anteriormente. El isomorfismo $v$ de $D^-(A)$ permite transportar $h_\nu$ a un endomorfismo de $K'$ en el sentido de $D^-(A)$. Se puede suponer que este último procede de un endomorfismo de complejos $h''_\nu:K'_\nu\to K'_\nu$. Apliquemos entonces el segundo lema con $L=K'_\nu$, $\varphi=h''_\nu$ y $\varphi_0=h'_{\nu-1}$; obtenemos la existencia de $h'_\nu:K'_\nu\to K'_\nu$ de modo que se satisfaga la condición 2); además, se satisfará la condición 3). Obsérvese que, con las notaciones del lema 1, si $M^i=0$ para un índice $i$, se deduce $L'^i=0$ gracias al lema de Nakayama; por tanto, todos los complejos $K'_\nu$ son efectivamente nulos fuera de un mismo intervalo de grados.
